\documentclass[11pt]{article}
\usepackage[T1]{fontenc}
\usepackage{lmodern,amsmath,amssymb,amsthm,microtype}
\usepackage[margin=1in]{geometry}
\usepackage{xcolor}
\definecolor{ink}{RGB}{24,64,84}
\usepackage[colorlinks=true,linkcolor=ink,urlcolor=ink,citecolor=ink]{hyperref}
\newtheorem{theorem}{Theorem}
\newtheorem{lemma}[theorem]{Lemma}
\newtheorem{corollary}[theorem]{Corollary}
\newcommand{\R}{\mathbb R}\newcommand{\Z}{\mathbb Z}
\newcommand{\E}{\mathbb E}\newcommand{\dd}{\,\mathrm d}
\newcommand{\pos}[1]{(#1)_+}
\setlength{\emergencystretch}{2em}
\title{The sharp support radius for an inverse-square sum of Fourier zeros}
\author{}\date{}
\begin{document}
\maketitle\vspace{-2em}
\begin{abstract}
Steinerberger asked whether a nonnegative function supported in a short arc can have a prescribed set of vanishing Fourier coefficients only if the radius dominates the reciprocal $\ell^2$ norm of the frequencies. In circle units the best universal constant is $1/(\pi\sqrt2)$. We give the sharp constant at each finite cardinality, characterize equality for probability measures, and bound the distance from a nearly extremizing measure to a pair of equal endpoint masses. The proof compares the ordered positive Fourier zeros with the odd multiples of the first zero of that pair. It applies to real frequencies and allows multiplicity.
\end{abstract}

\section{A support bound with its optimal constant}

Use $\mathbb T=\R/\Z$ and $\widehat\mu(a)=\int e^{2\pi iat}\dd\mu(t)$. Steinerberger's Problem~41 asks about $f\ge0$ of positive integral with support in $(-r,r)$ and prescribed zeros at a finite $A\subset\mathbb N$ \cite{Steinerberger}. We prove the proposed estimate for such functions by normalizing $f(t)\dd t$ to mass one. The closed interval formulation also describes the limiting measures that attain the sharp bound.

\begin{theorem}\label{thm:main}
Suppose $\mu$ is a probability measure on an interval of radius $r>0$ and $\widehat\mu$ has prescribed positive real zeros $a_1\le\cdots\le a_m$, each used at most its multiplicity. Then
\begin{equation}\label{eq:main}
 r\ \ge\ \frac{\bigl(\sum_{j=1}^m a_j^{-2}\bigr)^{1/2}}
 {4\bigl(\sum_{j=1}^m(2j-1)^{-2}\bigr)^{1/2}}
 \ >\ \frac{1}{\pi\sqrt2}\biggl(\sum_{j=1}^m a_j^{-2}\biggr)^{1/2}.
\end{equation}
Both the fixed-$m$ constant and the uniform constant are best possible, also for integer frequencies and smooth nonnegative functions in the sense of an infimum of radii. Equality in the first inequality holds precisely when
\[
 a_j=(2j-1)/(4r)\quad(1\le j\le m)
\]
and $\mu$ gives equal weight to the two endpoints of its containing interval.
\end{theorem}

The number $r$ refers to a containing interval, not the measure of an arbitrary support set. Translation multiplies the transform by a nonvanishing exponential. A measure on a circle arc therefore lifts to the stated real-line problem for integer frequencies. If an allowed radius on the circle exceeds $1/2$, one may first lift to a full fundamental interval.

The conclusion extends to every real $p>0$ at finite $m$:
\begin{equation}\label{eq:p}
 \sum_{j=1}^m a_j^{-p}\ \le\ (4r)^p\sum_{j=1}^m(2j-1)^{-p}.
\end{equation}
For $p>1$, the best uniform factor on the right, after division by $(4r)^p$, is $(1-2^{-p})\zeta(p)$. At $p\le1$ no finite uniform factor exists. One extremizing measure supplies all these sharp constants.

\section{From an integrated count to a comparison principle}

The analytic input in the ProveIt article is the integrated count of zeros for a positive sum of exponentials \cite{ProveIt}. Eremenko and Fryntov proved a related angular-distribution inequality for polynomials with nonnegative coefficients \cite{EF}. We reproduce the needed prefix bound to make the dependence explicit.

\begin{lemma}\label{lem:prefix}
Under the hypotheses of Theorem~\ref{thm:main},
\[
 \sum_{j=1}^k a_j\ \ge\ k^2/(4r)\qquad(1\le k\le m).
\]
\end{lemma}

\begin{proof}
For each prefix replace the measure by a finite positive measure on the same interval. Indeed, its values of the finitely many complex functionals $\int t^l e^{2\pi iat}\dd\mu(t)$ recording the required zero derivatives form a convex combination of point evaluation vectors, by compactness and Carath\'eodory's theorem. Translate its leftmost atom to zero. Translation preserves multiplicities. Write $0=t_0<\cdots<t_d=L\le2r$ for the nodes, $w_j>0$ for their masses and $G(s)=\sum_j w_j e^{t_j s}$. The zeros of $G$ include $2\pi i a_1,\dots,2\pi i a_k$ with the indicated multiplicities.

For $T>0$ one has the strip identity
\begin{equation}\label{eq:strip}
 0\le \int_{\R}\log\frac{G(x)}{|G(x+iT)|}\dd x
    = \frac{LT^2}{2}-2\pi\sum_{0<\operatorname{Im}\rho<T}(T-\operatorname{Im}\rho),
\end{equation}
where the sum is over all zeros of $G$ with multiplicity. To verify it, the two extreme terms dominate at the two ends of the real axis, uniformly in the imaginary coordinate. Thus the integral has exponentially convergent tails, and any logarithmic singularities are locally integrable. In distributions $\Delta\log|G|=2\pi\sum_\rho\delta_\rho$. Integrating the $xx$ derivative of $\log|G|$ over $x$ gives $L$, the difference of its limiting slopes. The second $T$ derivative of the left integral is consequently $L-2\pi\sum_\rho\delta_{\operatorname{Im}\rho}$ for $T>0$. There are no zeros in a neighborhood of the real axis, by the same uniform domination and positivity on the remaining compact interval. The integral and its first derivative vanish at $T=0$, by conjugation. Two integrations give the identity. Positivity of the integrand follows from the triangle inequality.

Put $\lambda_j=2\pi a_j$ and $\bar\lambda=k^{-1}\sum_{j\le k}\lambda_j$. The selected zeros contribute at least $\sum_{j\le k}(T-\lambda_j)_+\ge kT-k\bar\lambda$. Taking $T=2\bar\lambda$ in \eqref{eq:strip} gives $L\bar\lambda\ge\pi k$. Substituting $L\le2r$ proves the assertion.
\end{proof}

Normalize $x_j=4r a_j$ and $b_j=2j-1$. Lemma~\ref{lem:prefix} says exactly that the partial sums of the increasing sequence $(x_j)$ dominate those of $(b_j)$, since the sum of the first $k$ odd integers is $k^2$. Individual comparisons $x_j\ge b_j$ would be a stronger requirement than the proof needs. Define the functions
\[
 H_x(t)=\sum_{j=1}^m (t-x_j)_+,\qquad H_b(t)=\sum_{j=1}^m (t-b_j)_+ .
\]
Sorting both sequences gives the pointwise comparison
\begin{equation}\label{eq:hinge}
 H_x(t)=\max_{0\le k\le m}\biggl(kt-\sum_{j\le k}x_j\biggr)
 \ \le\ \max_{0\le k\le m}(kt-k^2)=H_b(t).
\end{equation}
This is the elementary form of weak supermajorization used here. For every $p>0$,
\[
 v^{-p}=p(p+1)\int_0^\infty(t-v)_+\,t^{-p-2}\dd t\quad(v>0).
\]
Integrating \eqref{eq:hinge} proves \eqref{eq:p}. The same argument applies to any nonnegative mixture of the functions $(t-v)_+$, including every twice continuously differentiable convex decreasing cost tending to zero at infinity by using its second derivative as the mixing density. The identity $\sum_{j\ge1}(2j-1)^{-2}=\pi^2/8$ gives \eqref{eq:main}.

\paragraph{Equality and sharpness.}
If equality holds in \eqref{eq:p}, the two continuous functions in \eqref{eq:hinge} coincide everywhere because the integrating density is strictly positive. The derivative jumps of these piecewise linear functions recover each node and its multiplicity, proving $x_j=b_j$. After centering the interval, the first zero at $a_1=1/(4r)$ gives $\int\cos(\pi t/(2r))\dd\mu(t)=0$. This cosine is nonnegative on $[-r,r]$, strictly positive in the interior. Thus all mass lies at the two endpoints. The sine moment forces the two masses to be equal. Conversely their transform is $\cos(2\pi r a)$, with simple positive zeros $(2j-1)/(4r)$. Taking $r=1/4$ gives all positive odd integers as zeros. Convolution with a nonnegative smooth probability density of arbitrarily small support preserves every one of these zeros and increases the enclosing radius by an arbitrarily small amount. This proves sharpness under the function and open-arc conventions. Letting $m\to\infty$ gives the asserted uniform statements.

\section{Stability of the extremizing measure}
Define the dimensionless finite-spectrum deficit
\[
 D=\sum_{j=1}^m\bigl((2j-1)^{-2}-(4r a_j)^{-2}\bigr).
\]
For probability measures on the line, $W_1$ denotes the infimum of $\mathbb E|X-Y|$ over all couplings. Center the support interval at zero and write $\nu_r=(\delta_{-r}+\delta_r)/2$.
\begin{theorem}\label{thm:stab}
If $D<1/2$, then
\[
 0\le4r a_1-1\le 2D,\qquad
 W_1(\mu,\nu_r)\le (2\pi+\pi^2/2)rD .
\]
The constants are uniform in the number of prescribed zeros.
\end{theorem}
\begin{proof}
Write $\delta=x_1-1\ge0$. For $1\le t\le 2$ and $k\ge2$, $kt-k^2\le0$. Hence $H_b=t-1$ there and $H_x=(t-1-\delta)_+$ by the prefix bound. With $d=\min(\delta,1)$,
\[
 D\ge 6\int_1^2\min(d,t-1)t^{-4}\dd t
       =1-(1+d)^{-2}-d/4\ \ge\ d/2.
\]
The last inequality follows by concavity on $[0,1]$ and the values at its endpoints. Thus $\delta\le2D$ whenever $D<1/2$.

Let $U$ have the law of $t/r$ under $\mu$. Since $\mathbb E e^{i\pi x_1 U/2}=0$, the Lipschitz bound for cosine gives $\mathbb E\cos(\pi U/2)\le(\pi/2)\delta$. On $[-1,1]$ this cosine dominates $1-|U|$ by concavity on either half. Write $d_0=\mathbb E(1-|U|)\le(\pi/2)\delta$. Send each $U$ to its sign $S$ with either endpoint chosen at zero. This costs $d_0$ in normalized units. The sine moment and the Lipschitz bound for sine give
\[
 |\mathbb E S|\le \mathbb E\big|S-\sin(\pi x_1 U/2)\big|
                 \le (\pi/2)(d_0+\delta).
\]
Moving the excess endpoint mass $|\mathbb E S|/2$ to the opposite endpoint costs $|\mathbb E S|$. Undoing the scale and using $\delta\le2D$ gives the claim.
\end{proof}

\section{Why dimension matters}
Steinerberger also suggests studying higher dimensions. The direct replacement of $a^{-2}$ by $|a|^{-2}$ in his inequality fails on $\mathbb T^d$ for every $d\ge2$. Indeed the two atoms at $\pm e_1/4$ have transform $\cos(\pi a_1/2)$, which vanishes at every lattice vector with odd first coordinate. Their reciprocal squared Euclidean norms have divergent sum already in the first two coordinate directions. For example the vectors with $2^k\le a_1<2^{k+1}$ odd and $2^k\le a_2<2^{k+1}$ contribute a positive absolute amount at each scale. A smooth convolution preserves all these zeros while keeping the support in a fixed ball of radius less than $1/2$. Thus the one-dimensional ordering in Lemma~\ref{lem:prefix} has genuine mathematical content. Determining a sharp higher-dimensional expression that takes the directions of the prescribed frequencies into account remains open here.

\begin{thebibliography}{9}
\bibitem{Steinerberger} S. Steinerberger, \emph{Open Problems}, 31 March 2026, Problem 41, pp.~41--42, \emph{A kind of uncertainty principle?}
\href{https://faculty.washington.edu/steinerb/openproblems.pdf}{Author's problem list}.
\bibitem{ProveIt} ProveIt Contributors, \emph{Beyond the Riemann hypothesis for Dirichlet $L$-functions}, revision \texttt{ff96fd44ab16}.
\href{https://github.com/VladimirReshetnikov/ProveIt/blob/ff96fd44ab160e65fb8a9d861243828d07c156cf/Analysis/GeneralizedRiemannHypothesis/Article/sections/fil.tex\#L223-L289}{Strip identity and mean-frequency theorem}.
\bibitem{EF} A. Eremenko and A. Fryntov, \emph{Remarks on the Obrechkoff inequality}, \href{https://arxiv.org/abs/1409.4849}{arXiv:1409.4849}.
\end{thebibliography}
\end{document}
